\scanpage{143}
\mspage{480}

Dans l'optique de la réflexion de ces derniers jours, c'est ce moment de retrouvailles avec mon enfance, crue perdue et morte une longue vie durant, qui marque la fin de la ``deuxième période'' de mon itinéraire spirituel: celle de la prédominance, dans ma vie personnelle, des m é c a n i s m e s é g o t i - q u e s, à l'encontre des forces créatrices, des forces de connaissance et de renouvellement, qui avaient passé par une stagnation presque complète de quarante ans.C'est aussi l'époque de la prépondérance d'une ``certaine force'', d'une force à caractère presque exclusivement ``viril'', à l'image des valeurs en \unsure{honneu|honneur} dans le monde environnant, aux dépens des aspects et forces profondes ``féminins'' de mon être, ignorés et refoulés (avec un succés jamais complet, Dieu merci!).

La toute première intuition sur la nature destructrice de cette force-là, laquelle avait dominé ma vie comme celle de ma mère, et celle d'autres femmes aussi qui avaient été importantes dans ma vie - cette intuition fait une courte apparition en ces jours de maturation intense, à la faveur sûrement du \unsure{ressurgis|ressurgis-}\newline sement de l'énergie yin, ``féminine'', dans mon appréhension consciente des \unsure{choses|choses.} Contrairement à ce dont je croyais hâtivement me souvenir tantôt, cette apparition n'a pas lieu dans la méditation de la veille des retrouvailles, mais quelques heures après celles-ci, dans une courte méditation sur le sens de ce qui venait de se passer. L'intuition naît et prend forme tout à la.fin des quelques pages de notes de cette méditation. Je perçois la nature destructrice de \unsure{ce|cet-}\newline te ``force'' (qu'aujourd'hui j'appellerais ``force superyang'', c'est à dire à dominante yang excessive) chez ma mère d'abord, puis chez d'autres femmes, pour \unsure{en|en-}chaîner avec ces lignes finales:

\begin{quote}
`` Quant à la ``force'' en moi-même, c'est elle certes qui a fait de moi la cible et l'objet, attendus une jeune vie durant, de la haîne secrète et du ressentiment de M., puis de J., puis de S. - d'une haîne déposée en elles longtemps avant qu'elles ne connaissent mon existence, en les jours désemparés d'une enfance privée d'amour.''
\end{quote}

Le mot ``enfance'', dans la dernière ligne encore qui témoigne d'un jour important entre tous dans ma vie, y apparaît d'ailleurs pour la dernière fois \unsure{pou|pour} près de trois ans! Quant à l'intuition sur la nature de la force superyang en moi, comme provocatrice de réactions antagonistes, voire de haîne et de ressentiment, elle a eu tendance (il me semble) à sombrer un peu dans l'oubli jusqu'à ces tout derniers jours encore. Plus précisément, elle est restée présente seulement dans ma perception de certaines relations importantes dans ma vie (et \unsure{su|sur-}\newline tout, des relations avec des femmes que j'ai aimées). Par contre, elle n'a \unsure{guèr|guère}

\mspage{142}
